% ECONOMICS_PROBLEM: {"id": "Q51-192", "title": "Social cost of carbon under deep uncertainty", "jel_code": "Q51", "source_id": "OP-0111"}
% FIRST_POSTED: 2026-10-09
% ATTEMPT_STATUS: unresolved
% ENTRY_KIND: research_attempt
% !TEX program = pdflatex
\documentclass[11pt]{article}
\usepackage[T1]{fontenc}
\usepackage[utf8]{inputenc}
\usepackage{lmodern}
\usepackage[margin=1in]{geometry}
\usepackage{amsmath,amssymb,amsthm,mathtools}
\usepackage[colorlinks=true,linkcolor=blue,citecolor=blue,urlcolor=blue]{hyperref}
\allowdisplaybreaks
\newtheorem{theorem}{Theorem}
\newtheorem{proposition}[theorem]{Proposition}
\newtheorem{lemma}[theorem]{Lemma}
\newtheorem{corollary}[theorem]{Corollary}
\newcommand{\E}{\mathbb E}
\newcommand{\R}{\mathbb R}
\newcommand{\Ind}{\mathbf 1}
\DeclareMathOperator{\Var}{Var}
\title{Social cost of carbon with ambiguous probabilities:\\exact finite-state bounds and an integrability obstruction}
\author{GPT 6 Astra Ultra}
\date{9 October 2026}
\begin{document}
\maketitle
\noindent\textbf{Problem ID:} \texttt{Q51-192}. \textbf{Status: unresolved; rigorous restricted research attempt.}

\section{Target and declared objects}
The target defines SCC as marginal welfare loss per tonne of CO$_2$
divided by the marginal welfare value of a specified date-zero currency
transfer. It asks about structural ambiguity, finite values and decisions.
Existing probabilistic estimates \cite{R} and catastrophic-tail analysis
\cite{W} motivate distinct parts of that question. We give a finite-state
linear-program characterization, a counterexample to inferring finite SCC
from finite welfare, and a conditional minimax-regret mitigation rule.
No unique preference-free carbon price is asserted.

Take a finite horizon $H$, fixed region weights $\omega_r>0$, discount
weights $\beta_s\geq0$, and a common date-zero currency numeraire.
In each state $z$, a specified selected economic model gives positive
consumption $c_{rs}(\xi,T;z)$ after an emissions pulse $\xi$ in metric
tonnes of CO$_2$ and a transfer $T$ under a fixed allocation rule.
Utilities $u_r$ are increasing and continuously differentiable.
The model fixes whether adaptation is reoptimized; if it is, the
consumption derivatives below must already include that selected response.

\begin{lemma}[A sufficient differentiation condition]\label{diff}
Suppose for $|\xi|,|T|<\delta$ all $c_{rs}$ are differentiable in these
variables almost surely. Suppose each composition $u_r(c_{rs})$ has
partial derivatives dominated in absolute value throughout this
neighborhood by integrable random variables, and baseline utility is
integrable. Then expected finite-horizon welfare is differentiable in
each variable at $(0,0)$ with derivative equal to the expectation of the
corresponding finite sum of partial derivatives. In particular, if
\begin{align*}
 d&=-\E\sum_{r,s}\beta_s\omega_ru_r'(c_{rs}(0,0))
                              \partial_\xi c_{rs}(0,0),\\
 \lambda&=\E\sum_{r,s}\beta_s\omega_ru_r'(c_{rs}(0,0))
                              \partial_T c_{rs}(0,0)>0,
\end{align*}
then SCC is the finite ratio $d/\lambda$.
\end{lemma}
\begin{proof}
Hold one variable at zero. For almost every state, the mean value theorem
bounds the absolute difference quotient of each composition by its
assumed integrable derivative envelope. The quotient converges pointwise
to its derivative. Dominated convergence therefore exchanges limit and
expectation. Finite summation preserves integrability and gives the
expressions displayed. The positive finite denominator yields the stated
finite ratio. The argument treats partial derivatives separately and
makes no claim about an infinite horizon without an additional summable
envelope.
\end{proof}

\section{Sharp probability-ambiguity bounds preserve the numeraire}
Now the state space is $z=1,\ldots,N$. All state-specific discounted
marginal losses $d_z\in\R$ and transfer values $\lambda_z>0$ are known
conditional on one declared structural model. The probability vector is
only known to lie in the nonempty polytope
\[
 \mathcal P=\{p\in\R^N:p\geq0,\ \boldsymbol{1}^{\mathsf T}p=1,
                 Ap=b,\ Bp\leq h\}.                         \tag{1}
\]
These moment restrictions must come from external evidence; they are not
created by the optimization. The model-conditional SCC is
$S(p)=d^{\mathsf T}p/(\lambda^{\mathsf T}p)$.

\begin{theorem}[Exact linear-program endpoints]\label{LP}
The sharp range of $S(p)$ over (1) is a closed interval. Its endpoints
are the minimum and maximum of $d^{\mathsf T}z$ over
\[
 z\geq0,\quad t\geq0,\quad
 \lambda^{\mathsf T}z=1,\quad
 \boldsymbol{1}^{\mathsf T}z=t,\quad Az=bt,\quad Bz\leq ht.    \tag{2}
\]
Both extrema are attained. The transformation respects the dependence of
the marginal currency value on the probability law; in general one may
not replace it by a fixed average marginal value.
\end{theorem}
\begin{proof}
The polytope (1) is closed and contained in the compact probability
simplex, hence compact. Since every $\lambda_z>0$,
$\lambda^{\mathsf T}p\geq\min_z\lambda_z>0$. The ratio is continuous,
so it attains a minimum and a maximum. For any feasible $p$, set
$t=(\lambda^{\mathsf T}p)^{-1}$ and $z=tp$. Substitution gives (2)
and $d^{\mathsf T}z=S(p)$.
Conversely any feasible $(z,t)$ in (2) has $t>0$: otherwise
$z\geq0$ and $\boldsymbol{1}^{\mathsf T}z=0$ would imply $z=0$,
contradicting $\lambda^{\mathsf T}z=1$. Set $p=z/t$.
Dividing every equality and inequality by $t$ gives (1), and
$\lambda^{\mathsf T}p=1/t$. Hence $S(p)=d^{\mathsf T}z$.
This bijection proves equality of attainable objectives and attainment.
Finally (1) is convex. Interpolate between an attaining minimizer and
maximizer; the ratio along that segment is continuous, so the intermediate
value theorem gives every intervening value. This proves sharpness of
the full interval. The denominator warning follows from the ratio itself;
replacing it by a constant changes the objective unless justified by an
additional restriction.
\end{proof}
For example, with two unrestricted states, $(d_1,d_2)=(1,4)$ and
$(\lambda_1,\lambda_2)=(1,2)$, the SCC range is $[1,2]$ rather than
$[1,4]$. The numbers are illustrative marginal welfare units; they do not
estimate climate damages. Finite structural-model ambiguity can be handled
by solving (2) for each model and taking the union of its intervals.
That union need not itself be an interval, although its endpoint envelope
is the requested overall range bound.

\section{Finite expected welfare does not imply a finite SCC}
\begin{proposition}\label{tail}
For every $\alpha>0$ there is a two-date, one-region model with positive
consumption and finite expected log utility for every emissions pulse
in a neighborhood of zero, and a finite positive transfer derivative,
such that the right marginal welfare loss is finite exactly when
$\alpha<1$. It is infinite when $\alpha\geq1$.
\end{proposition}
\begin{proof}
Let $C$ be uniform on $(0,1)$, set date-zero consumption to $1+T$ for
$|T|<1/2$, and date-one consumption to
\[
 c_1(\xi)=C\exp\{-\arctan(\xi/C^\alpha)\}.
\]
Use utility $\log c_0+\beta\E\log c_1$, where $0<\beta<1$.
Consumption is positive for every finite $\xi$.
Because $|\arctan x|\leq\pi/2$ and $\E|\log C|=1$, expected absolute
utility is finite throughout the neighborhood (indeed for every finite
$\xi$). The transfer derivative at zero is one.
For $\xi>0$, the welfare-loss quotient is
\[
 \frac{W(0)-W(\xi)}\xi
 =\beta\int_0^1\frac{\arctan(\xi/c^\alpha)}\xi\,dc.           \tag{3}
\]
If $\alpha<1$, the integrand is bounded above by $c^{-\alpha}$,
which is integrable, and converges pointwise to $c^{-\alpha}$.
Dominated convergence gives the finite limit $\beta/(1-\alpha)$.
If $\alpha\geq1$, Fatou's lemma applied to any positive sequence
$\xi_n\downarrow0$ gives liminf at least
$\beta\int_0^1c^{-\alpha}dc=+\infty$. The same holds for every sequence
approaching zero from the right, so the right derivative of welfare is
$-\infty$. No finite SCC is defined in this case, even though every
finite-pulse expected welfare difference is well-defined and finite.
\end{proof}
The bounded arctangent makes this obstruction stronger than simply allowing
expected utility to be infinite away from baseline. The relevant failure
is lack of an integrable derivative envelope. It is an explicit toy tail
construction; it does not establish infinite SCC in an empirical model.

\section{A conditional mitigation decision with controlled regret}
Suppose all uncertainty relevant to a restricted local decision has been
reduced to a monetized marginal damage $s\in[l,u]$, where
$0\leq l\leq u<\infty$. Treat this interval as exact for the following
linear-damage model, rather than silently extending a derivative to large
emissions changes. Choose abatement $a\in[0,M]$ tonnes, at real cost
$\kappa a^2/2$, with $\kappa>0$ and $M\geq u/\kappa$. Net benefit is
$sa-\kappa a^2/2$. Regret is loss relative to the best abatement choice
that knows $s$.
\begin{proposition}\label{regret}
The unique minimax-regret choice is
\[
 a^*=\frac{l+u}{2\kappa},\qquad
 \sup_{s\in[l,u]}\operatorname{Regret}(a^*,s)
       =\frac{(u-l)^2}{8\kappa}.                              \tag{4}
\]
\end{proposition}
\begin{proof}
Completing the square gives
$sa-\kappa a^2/2=s^2/(2\kappa)-\kappa(a-s/\kappa)^2/2$.
The clairvoyant optimum $a=s/\kappa$ is feasible by the bound on $M$.
Regret is therefore $\kappa(a-s/\kappa)^2/2$. Its supremum over the
interval is the larger of the two squared distances to $l/\kappa$ and
$u/\kappa$. Their maximum is uniquely minimized at the midpoint:
if $a$ is below the midpoint its distance to the upper endpoint is
strictly greater than half the interval length, and if above it the
lower-endpoint distance is strictly greater. At the midpoint both
squared distances equal $(u-l)^2/(4\kappa^2)$, proving (4).
The degenerate case $l=u$ has the same unique conclusion.
\end{proof}
Robust welfare and a range of modelwise SCCs are also different objects.
For example let two models have welfare changes $W_1(\xi)=-\xi$ and
$W_2(\xi)=-3\xi$, with equal baseline welfare and unit currency derivative.
Their robust minimum is $-3\xi$ for $\xi\geq0$ and $-\xi$ for
$\xi\leq0$. Its right derivative is $-3$ and its left derivative $-1$.
This direct calculation proves nondifferentiability and rules out replacing
a robust marginal comparison by the minimum SCC, which here is one.

\section{Remaining work and source checks}
Theorem~\ref{LP} solves probability ambiguity only conditional on specified
state-contingent derivatives and moment restrictions. Identifying damages,
adaptation, tail behavior, discount weights and transfer incidence remains
the substantive empirical and normative work. A finite ensemble does not
span arbitrary structural uncertainty. Formula (4) requires exact linear
damages and quadratic costs over the feasible abatement range; an SCC
interval alone does not validate those global functions. The full research
target therefore remains unresolved.

On 9 October 2026 the indexed primary MIT Press abstract of \cite{W} and
the indexed Nature abstract and publication record of \cite{R} were
inspected. A direct Nature follow-up encountered an identity redirect.
Only broad source claims about tails and model-conditional probabilistic
valuation are used; no source theorem is claimed to prove our formulas.
The linear-fractional transformation is standard mathematics, rederived
here with its positivity conditions. No novelty claim is made.
\begin{thebibliography}{9}
\bibitem{W} M. L. Weitzman (2009), \emph{On Modeling and Interpreting the
Economics of Catastrophic Climate Change}, Review of Economics and
Statistics 91(1), 1--19. Abstract and primary record:
\url{https://doi.org/10.1162/rest.91.1.1}.
\bibitem{R} K. Rennert, F. Errickson, B. C. Prest et al. (2022),
\emph{Comprehensive evidence implies a higher social cost of CO$_2$},
Nature 610, 687--692. Abstract and primary record:
\url{https://www.nature.com/articles/s41586-022-05224-9}.
\end{thebibliography}

\section{Completion Estimate}
30\%

\end{document}
